% Chapter 1

\chapter{Estado del arte} % Main chapter title

\section{Revisión Bibliográfica y Estado del Arte}


Esta discrepancia entre ruta ejecutada no es exclusiva de SimpliRoute, sino que ha sido documentada de forma sistemática en la última milla bajo distintos nombres, como implicit preferences\cite{canoy2023learn}, tacit knowledge \cite{merchan20242021} o context-aware preferences \cite{cook2024constrained}, a partir del desafío de datos de Amazon, cuantificaron que las rutas ejecutadas por conductores reales son en promedio un 12,9 \% más largas en tiempo de viaje que la solución óptima teórica, evidencia de que el conductor no minimiza linealmente el tiempo de trayecto sino que opera bajo criterios implícitos. La magnitud del fenómeno en SimpliRoute (85,81 \% de desviación) sugiere que, lejos de ser una anomalía marginal, el motor de ruteo está resolviendo sistemáticamente un problema distinto al que el conductor enfrenta en terreno.

La literatura científica contemporánea en logística de última milla ha transitado desde modelos prescriptivos hacia marcos adaptativos basados en datos. Así, estos enfoques modernos buscan capturar la incertidumbre del entorno urbano, las dinámicas operacionales complejas y el conocimiento tácito humano a través de un historial de datos pasados. De manera que, para fundamentar metodológicamente esta investigación, el estado del arte se divide en tres pilares esenciales. Primero, el aprendizaje

\section{Literatura}


Ahora bien, para resolver este problema en el contexto de Simpliroute se requiere considerar tres aspectos clave. Primero, heterogeneidad individual, es decir, la diferencia en el comportamiento de los conductores frente a la misma trayectoría no es la misma, donde idealmente se pueda aprender el perfil único de cada conductor, no un promedio de la zona. Segundo, lograr una escalabilidad, debe funcionar con una gran cantidad de carteras de clientes que cambian todos los días y trazas GPS dispersas. Y, por último, reinyectabilidad operacional, donde el resultado sea directamente consumible por el motor de optimización de Simpliroute.
\subsection{Aprendizaje en Línea de Tiempos de Viaje mediante Bandidos Multibrazo Combinatorios}
 
Una vertiente fundamental de la literatura aborda el problema desde la perspectiva
del aprendizaje adaptativo de la red de transporte, en escenarios donde el operador
logístico posee información limitada o estocástica sobre los tiempos reales de
viaje \cite{lagos2024online}. Es decir, en términos simples, un bandido
multibrazo es un modelo de decisión secuencial en el que un agente tiene dos fases, una de exploración y otra de explotación \cite{robbins1952some}. Por un lado, en la fase de exploración, el agente debe elegir repetidamente entre distintas alternativas, sin conocer de antemano cuál es la mejor, aprendiendo únicamente a partir de la retroalimentación que recibe después de cada elección. Por otro lado, en la fase de explotación, el agente utiliza la información aprendida para tomar decisiones óptimas. 
 
Así, \citet{lagos2024online} formalizan el problema del camino más corto modelándolo como un problema de bandidos multibrazo combinatorios (CMAB)
bajo un entorno de retroalimentación de tipo semibandit. En este caso, cada arco de la red vial actúa como un brazo cuyo tiempo de viaje real solo se revela después de haber sido recorrido. Donde, la metodología  de exploración y explotación son basadas en el algoritmo de Thompson sampling.

\subsection{Enfoques Predictivos-Prescriptivos y Aprendizaje de Preferencias Multi-Objetivo}
 
Una línea de investigación complementaria se enfoca en el aprendizaje de
preferencias a nivel macro, es decir, a nivel del conjunto completo de decisiones
de una plataforma logística en lugar de a nivel de un conductor individual.
Nourmohammadi, Hu, Rey y Saberi (2025) introducen un enfoque de optimización guiado
por datos para inferir implícitamente las ponderaciones de las funciones objetivo
en un Problema de Ruteo Vehicular Capacitado con Ventanas de Tiempo Multi-Objetivo
(MO-CVRPTW). A diferencia de los métodos de Optimización Inversa (IO) —que se
detallan en la sección~\ref{subsec:optimizacion-inversa} y que suelen operar sin
una interpretación clara de los pesos aprendidos—, este marco utiliza
escalarización y normalización de objetivos para resolver la disparidad de escalas
numéricas entre las distintas métricas de las funciones objetivo.
 
Tras evaluar estadísticamente un conjunto de escenarios históricos, los autores
demostraron que las preferencias de los tomadores de decisiones logísticas se
concentran en tres objetivos principales: la minimización de la distancia total de
viaje, la minimización del tamaño de la flota utilizada y el balance de la carga de
trabajo. Posteriormente, utilizan el algoritmo de bosques aleatorios (RF) para
predecir las funciones objetivo a partir de las características intrínsecas de las
instancias de entrega, eliminando la necesidad de recompilar el motor de
optimización en cada exploración de pesos. Una premisa fundamental de este enfoque,
sin embargo, es asumir homogeneidad en las preferencias de la flota: el modelo no
busca parametrizar el comportamiento aislado de cada conductor, sino capturar la
visión unificada del tomador de decisiones.
 

------------------------------------------------------------
\subsection{Búsqueda Local Restringida Basada en Penalizaciones Jerárquicas}
 
Un hito fundamental en el modelamiento del comportamiento del conductor a nivel
micro —es decir, en la secuenciación detallada de paradas dentro de una ruta ya
asignada— fue establecido por Cook, Held y Helsgaun (2024), quienes desarrollaron
el algoritmo LKH-AMZ para competir en el Amazon Last Mile Routing Research
Challenge de 2021 (Merchán y cols., 2024). Los autores abordan el problema de la
secuencia detallada de un único vehículo tratándolo formalmente como un Problema
del Vendedor Viajero Asimétrico (ATSP, por su sigla en inglés): una variante del
clásico problema del vendedor viajero en la que el costo de ir de un punto $A$ a un
punto $B$ no necesariamente es igual al costo de ir de $B$ a $A$, lo cual refleja
de forma más realista las restricciones de tránsito urbano (calles de un solo
sentido, giros prohibidos, congestión asimétrica).
 
Para aproximar las soluciones teóricas a las decisiones reales, los autores
propusieron un enfoque de optimización local restringida por penalizaciones: en
lugar de minimizar exclusivamente la distancia o el tiempo de viaje, el algoritmo
penaliza las secuencias que se alejan de los patrones de zonificación geográfica
observados históricamente en los datos de Amazon. El algoritmo LKH-AMZ obtuvo el
primer lugar global de la competencia, alcanzando un \textit{score} de 0.0248
frente a los datos ocultos de Amazon, y demostró además una alta eficiencia
computacional al procesar instancias complejas en fracciones de segundo.
  
\subsection{Aprendizaje de Preferencias Implícitas mediante Modelos Probabilísticos de Transición}
En la práctica logística, un planificador experimentado rara vez acepta sin modificaciones la ruta que propone un software de optimización, la ajusta según su conocimiento del terreno, la historia de cada cliente y las restricciones tácitas que ha aprendido con el tiempo. Un sistema que pueda aprender esas modificaciones de forma automática, sin que el planificador deba articular explícitamente sus razones, representaría una mejora significativa para la adherencia operacional que buscamos resolver.

Con ese objetivo, \cite{canoy2023learn} proponen un marco de aprendizaje de preferencias implícitas para el VRP basado en un modelo de transición markoviano entre paradas. Donde, la idea central consiste en reemplazar la matriz de distancias del solver VRP por una estimada a partir de las rutas manuales históricas. cuanto más frecuentemente un planificador haya encadenado la parada entonces será mayor la probabilidad de transición a esa zona de la parada. 

Ahora bien, esta matriz aprendida se integra directamente con la infraestructura de optimización existente, combinando el criterio objetivo de distancia con el criterio subjetivo de preferencia, sin necesidad de reformular matemáticamente el VRP subyacente. En su validación sobre datos reales de una empresa de transporte, los autores demuestran que las rutas generadas con la matriz aprendida se acercan sistemáticamente más a las soluciones manuales del planificador que las rutas puramente basadas en distancias, incluso ante cambios en el conjunto de clientes.

No obstante, este enfoque presenta una limitación fundamental en el contexto de SimpliRoute. Ya que, el modelo captura la ocurrencia histórica entre arcos, pero no la causa subyacente de cada preferencia. De manera que, si dos conductores distintos han encadenado con frecuencia las mismas paradas por razones completamente diferentes, entonces el modelo agrega ambos patrones sin distinguirlos. Por lo tanto, esta mezcla de motivos heterogéneos deteriora la calidad de la predicción cuando distintos conductores operan sobre la misma cartera de clientes.


\subsection{Optimización Inversa Aplicada al Ruteo de Vehículos}

Cuando un conductor experto elige una ruta, implícitamente está resolviendo un problema de optimización cuya función objetivo no está siendo modelada matemáticamnte. La Optimización Inversa (IO) invierte esta lógica, es decir, en lugar de resolver un problema de optimización dado un objetivo conocido, infiere el objetivo a partir de las soluciones observadas \cite{ahuja2001inverse}. De manera que, aplicada al ruteo, permite estimar qué pesos latentes sobre distintos atributos viales y operacionales habrían producido las rutas ejecutadas históricamente por un conductor.

Ahora bien, esta metodología ha sido formalizada para el VRP \cite{zattoni2025inverse}, proponiendo una metodología de IO con función hipótesis, función de pérdida y algoritmo de optimización estocástico de primer orden adaptados a las restricciones estructurales del problema de ruteo. El marco parte del supuesto de que las rutas históricas ejecutadas por los conductores son soluciones aproximadamente óptimas de un VRP cuya función de costo verdadera es desconocida. Así, a través de IO, se estima el vector de pesos latente que refleja las preferencias implícitas del decisor. En su validación sobre el desafío de última milla de Amazon (2021), los autores aprenden las preferencias de los conductores en términos de zonas geográficas predefinidas de la ciudad, alcanzando el segundo lugar entre los 48 equipos finalistas con un puntaje de 0.0302.

Sin embargo, existen dos limitaciones relevantes frente al contexto de SimpliRoute. Primero, la implementación principal del marco depende de la existencia de una zonificación geográfica preestablecida. Ya que, en SimpliRoute, las bodegas de los clientes y los manifiestos de carga no cuentan con una estructura de zonas previamente definida, equivalente a la de Amazon, lo que impide la transferencia directa del enfoque. Segundo, la IO aprende un único vector de parámetros para el conjunto de todas las rutas analizadas, representando un promedio de la flota más que el perfil individual de cada conductor. Por lo tanto, tanto la dependencia de la zonificación, como la homogeneización de la flota, constituyen brechas directas respecto de los requisitos del problema bajo estudio.

\subsection{Aprendizaje por Refuerzo Inverso y Aprendizaje Profundo Basado en Preferencias con Trazas GPS}
El Aprendizaje por Refuerzo Inverso (IRL) parte de una premisa intuitiva, si se observa el comportamiento de un agente que toma decisiones de forma experta, es posible inferir la función de recompensa que dicho agente está maximizando, aunque esa función nunca haya sido formulada explícitamente \cite{ng2000algorithms}. 

En el dominio del ruteo vehicular, esto significa que a partir de un historial de trazas GPS es posible aprender qué combinación latente de factores viales está pesando el conductor al elegir sus rutas, sin que él deba declararlo.

Sin embargo, aun presenta ciertas complicaciones, donde identifican dos limitaciones técnicas que afectan a los métodos IRL estándar cuando se aplican sobre trazas GPS reales en redes viales urbanas \cite{zhang2024recursive}. Por un lado, estimar las frecuencias esperadas de visita a cada segmento vial, requiere aproximaciones numéricas iterativas que son simultáneamente costosas en cómputo e imprecisas. Por otro lado, la cobertura de las trazas GPS está sesgada hacia los tramos más transitados, dificultando el aprendizaje de preferencias en segmentos poco frecuentados.

Así, para superar ambas limitaciones, los autores proponen combinar el modelo de logit recursivo, que deriva analíticamente las frecuencias esperadas, reduciendo el tiempo de cómputo en más de un 95\%, el cual transfiere el conocimiento adquirido en tramos populares hacia tramos escasamente cubiertos. De manera que, la validación sobre datos GPS reales de la ciudad de Chengdu muestra que las rutas planificadas con este enfoque presentan mayor correspondencia con las trayectorias ejecutadas que los métodos de IRL de referencia.

Por su parte, \cite{pan2025preference} proponen un enfoque de Aprendizaje por Refuerzo Profundo Basado en Preferencias (PbRL), cuya función de recompensa está especializada para aprender de comparaciones entre rutas históricas completas. Así, el modelo utiliza una red neuronal entrenada sobre el Problema de Ruteo Vehicular Capacitado (CVRP). Donde, demuestra que incorporar señales de preferencia humana permite generar soluciones más alineadas con el comportamiento histórico de los conductores que los métodos de aprendizaje para optimización, los cuales se limitan a minimizar distancia o tiempo de viaje.

No obstante, ambos enfoques presentan limitaciones críticas frente al problema de SimpliRoute. El método de \cite{zhang2024recursive} fue concebido y validado para ruteo de un único par origen-destino en redes urbanas (conducción diaria y \textit{ride-hailing}), un contexto donde la decisión consiste en elegir un camino entre dos puntos fijos sobre un grafo vial. Pero, el problema de SimpliRoute tiene una estructura combinatoria distinta: no se trata de elegir un camino, sino de determinar la secuencia óptima entre múltiples paradas (VRP), un espacio de decisión exponencialmente mayor para el cual estos modelos no han sido adaptados. En cuanto al enfoque de \cite{pan2025preference}, la función de recompensa aprendida compara rutas completas entre sí sin descomponerla en los factores viales individuales que la gobiernan, lo que limita la reinyectabilidad del modelo directamente en la heuristica de Simpliroute.



\section{Brecha de Conocimiento}

A partir de la revisión de las cinco familias metodológicas presentadas, se constatan tres limitaciones estructurales que ninguna de ellas resuelve de forma simultánea frente al problema específico de SimpliRoute.

Primera limitación: heterogeneidad individual no modelada.
Los métodos de aprendizaje de preferencias multiobjetivo a nivel de plataforma (Nourmohammadi y cols., 2025) y la IO aplicada a ruteo (Zattoni Scroccaro y cols., 2025) estiman un conjunto único de parámetros para toda la flota, asumiendo homogeneidad entre conductores. Sin embargo, el análisis exploratorio desarrollado en esta tesis evidencia que distintos operadores que cubren una misma zona de reparto exhiben patrones de secuenciación sistemáticamente diferentes: la variación en el uso de tipos de vía no es aleatoria sino consistente por conductor. Un modelo que promedia la flota introduce ruido y oculta los perfiles individuales que determinan la adherencia real.

Segunda limitación: imposibilidad de separar el efecto de zona del efecto de preferencia personal.
Los métodos jerárquicos basados en zonificación (Cook y cols., 2024) y la IO con zonas preestablecidas (Zattoni Scroccaro y cols., 2025) descansan en una estructura geográfica predefinida que no existe en SimpliRoute. Más relevante aún, el análisis exploratorio de esta tesis identifica que parte de la variación observada en el comportamiento de los conductores es atribuible a la asignación de zona —un efecto estructural y colectivo— y no exclusivamente a la preferencia personal. Ninguno de los métodos revisados ofrece un mecanismo para disociar ambos efectos, lo cual puede llevar a atribuir patrones geográficos de la zona a preferencias personales del conductor, o viceversa.

Tercera limitación: ausencia de reinyectabilidad operacional en el motor VRP.
Los enfoques de IRL y PbRL (Zhang y cols., 2024; Pan y cols., 2025) aprenden una representación de la preferencia del conductor a partir de trazas GPS, pero dicha representación no está formulada como una corrección directamente aplicable a la función de costo de un solver VRP existente. En el contexto operativo de SimpliRoute, el resultado metodológico debe ser consumible por el motor de optimización —ya sea como un vector de pesos ajustado en la función de costo o como una restricción blanda aprendida— sin requerir el rediseño completo del sistema de ruteo.

En síntesis, existe una brecha metodológica no cubierta en la intersección de estas tres condiciones: un marco capaz de inferir preferencias latentes \textbf{a nivel individual por conductor}, \textbf{distinguiendo el efecto estructural de la asignación de zona del efecto personal}, y cuyo resultado sea \textbf{directamente reinyectable en un motor de optimización VRP}. Esta investigación se posiciona en esa intersección, aprovechando además la disponibilidad infrecuente en la literatura de trazas GPS de alta frecuencia con contrapartes de rutas planificadas por conductor —un tipo de dato señalado repetidamente como la principal barrera para avanzar en esta línea (Merchán y cols., 2024; Cook y cols., 2024).